
\usepackage{xeCJK}%字體設定
\setCJKmainfont{cwTeXKai}[AutoFakeBold=1.5]


\setmainfont{Times New Roman}  % 英文正文字體

\usepackage[top=2cm,bottom=2cm,left=2cm,right=2cm,a4paper]{geometry}%版面設定

\usepackage{indentfirst}%縮排

\usepackage{setspace}%行距設計
\onehalfspacing%半行距
\setlength{\parskip}{3pt}%段落之間

\usepackage{moresize}%字體大小

\usepackage{graphicx}

\usepackage{amssymb}%數學符號

\usepackage{amsmath} % 數學符號
\usepackage{booktabs}         % 漂亮的表格（\toprule, \midrule, \bottomrule）
\usepackage{siunitx}          % 國際單位制 (建議處理單位更漂亮)

\usepackage{multicol}%引入函數使下面那些可以變成雙欄

\usepackage{float}%圖片固定

\usepackage{subcaption}

\usepackage{zhnumber} % 引入中文數字套件

\usepackage{tcolorbox}



% --- 程式碼樣式 ---
% 設定 listings 用的字體
\usepackage{listings}          % 插入程式碼
\usepackage{xcolor}            % 程式碼顏色
\setmonofont{Courier New} % 設定等寬字體
\lstset{
  basicstyle=\ttfamily\small, % 使用等寬字體顯示程式碼
  numbers=left,
  numberstyle=\tiny,
  keywordstyle=\color{blue},
  commentstyle=\color{green!50!black},
  stringstyle=\color{red},
  breaklines=true,
  frame=single,
  tabsize=2,
  showspaces=false,      % 不要把空格顯示成 ␣
  showstringspaces=false % 不要在字串裡顯示 ␣
}

% \renewcommand{\thesection}{\Roman{section}}

\usepackage{fancyhdr}
\usepackage{hyperref}
\usepackage{xcolor}

\definecolor{os_blue}{HTML}{1B263B} % 你可以根據選定的配色調整

\pagestyle{fancy}
\fancyhf{}
\renewcommand{\headrulewidth}{0pt}
\fancyfoot[L]{\footnotesize \textcolor{gray}{© 2026 \textbf{charlie.note.os}}}
\fancyfoot[C]{\footnotesize \textcolor{gray}{CC BY-NC-SA 4.0}}
\fancyfoot[R]{\footnotesize \textcolor{gray}{Page \thepage}}

% 在講義末頁可以加入超連結
\newcommand{\copyrightnotice}{
    \vfill
    \begin{center}
        \scriptsize \textcolor{gray}{
            本作品採用創用 CC 姓名標示-非商業性-相同方式分享 4.0 國際授權釋出。\\
            詳情請見：\href{https://creativecommons.org/licenses/by-nc-sa/4.0/deed.zh_TW}{CC BY-NC-SA 4.0 條款說明}
        }
    \end{center}
}



\renewcommand{\thesection}{\zhnumber{\arabic{section}}、}
\renewcommand{\thesubsection}{(\arabic{subsection})}
\renewcommand{\thesubsubsection}{\arabic{subsubsection}.}




\newtcolorbox{concept}{
colback=green!5,
colframe=green!50!black
}

\newtcolorbox{trick}{
colback=yellow!10,
colframe=orange!60
}